Goiás and the Federal District ($\sim$340,000~km\textsuperscript{2} of the Cerrado biome)
concentrate the central tension of land-use planning in the tropics: expanding and intensifying
production without compromising the ecosystems and water resources it depends on. This
dissertation builds, for GO~and~DF at 250~m, an \textbf{open multi-segment agricultural
suitability atlas with zoning and climate projection}, covering seven land-use segments:
soybean, sugarcane, other annual crops, aquaculture, livestock, forest/Cerrado
conservation, and photovoltaic generation. Current suitability results from an FAO-style
multicriteria evaluation (fuzzy membership and AHP weights). Agro-environmental zoning comes from
unsupervised k-means clustering over a decorrelated representation of the variables. And the
mid-century projection applies the change-factor method to an ensemble of five
CMIP6 models. All processing runs server-side, on Google Earth Engine, \textbf{with no
upload of external data}, validated against current land use (MapBiomas) and a
productivity indicator (MODIS MOD17). The results answer the three research questions: current
suitability surfaces for the seven segments and a partition into \textbf{seven zones} with
distinct comparative vocations, among them a high plateau whose best vocation is conservation. The
zones, built solely from biophysical variables, reproduce the observed geography of land use
without using land cover as an input (soybean AUC of 0.83). They also reveal a relevant gap
between potential suitability and current use, especially in pastureland, pointing to
intensification --- rather than expansion onto native vegetation --- as the most likely direction
of land-use change in the region. Mid-century climate
change is modest,
yet directionally robust (inter-model agreement $\geq$ 0.92): it affects mainly
conservation and, to a lesser degree, the maize/cotton/sorghum group and photovoltaic generation,
while the prime soybean and sugarcane plateaus remain essentially stable. The method runs entirely
in the cloud from public data, with no upload, making it reproducible at low marginal cost for
other Cerrado states.

\vspace{1em}
\noindent\textbf{Keywords:} agricultural land suitability; land evaluation (FAO);
analytic hierarchy process (AHP); agro-environmental zoning; Cerrado; Google Earth Engine;
CMIP6 climate projection; multicriteria decision analysis.
